\documentclass[12pt]{article}
\usepackage[utf8]{inputenc}
\usepackage[T1]{fontenc}
\usepackage{hyperref}
\title{Tiranige Diga Dogon Language Resources for K-12 Curriculum Development}
\author{Compiled by Claude Sonnet 5 (claude-sonnet-5) for Andr\'as Kornai}
\date{}
\begin{document}
\maketitle
\section{Introduction}
Tiranige Diga Dogon (also called Duleri; ISO 639-3 tde, Glottocode tira1258) is a western Dogon language of the Bandiagara Plateau in Mali. One source reports about 4,200 speakers; Joshua Project gives 9,700 for the people group. This report is based on one web search.
\section{Primary Educational Resources}
No primers, textbooks or school materials were found.
\section{Linguistic and Reference Materials}
\begin{itemize}
\item Jeffrey Heath, \emph{A grammar of Tiranige (Dogon, Mali)}, University of Michigan, open access (CC BY 4.0). The page also points to a lexicon and audio files: \url{https://deepblue.lib.umich.edu/items/d11bf892-b231-48c0-b264-0cd396acebcd}
\item Numeral system description, MPI-EVA: \url{https://lingweb.eva.mpg.de/channumerals/Dogon-Tiranige-Diga.htm}
\item Wikipedia, Duleri Dogon: \url{https://en.wikipedia.org/wiki/Duleri_Dogon}
\item Zenodo record (title not shown in the results): \url{https://zenodo.org/record/6350672}
\end{itemize}
\section{Religious and Community Resources}
Joshua Project reports that Scripture translation has started: \url{https://www.joshuaproject.net/index.php/people_groups/20583}. No scripture text was seen.
\section{Digital and Audio Resources}
Heath's audio files and lexicon are said to be in Deep Blue (see above). The Endangered Languages Project has an entry: \url{https://endangeredlanguages.com/node/95590}
\section{Curriculum Development Context}
The language has a full descriptive grammar but no evidence of a standard orthography or school materials. Any curriculum would rest on Heath's grammar, lexicon and audio.
\end{document}